\documentclass[../../main.tex]{subfiles}% 注意这里的文件路径不能用 ./main.tex ,否则用latexmk编译子文件会报错
\graphicspath{{\subfix{./image/}}} % 指定图片目录,后续可以直接使用图片文件名
% 注意这里的文件路径不能用 ../../image/ ,否则用latexmk编译子文件会报错

% 例如:
% \begin{figure}[H]
% \centering
% \includegraphics[scale=0.3]{图.png}
% \caption{}
% \label{figure:图}
% \end{figure}
% 注意:上述\label{}一定要放在\caption{}之后,否则引用图片序号会只会显示??.

\begin{document}

\section{$L^p$空间的定义与不等式}

\begin{definition}[$L^p$空间]
(i) 设 $f(x)$ 是 $E\subset\mathbb{R}^n$ 上的可测函数，记
\begin{align*}
\|f\|_p = \left(\int_E |f(x)|^p \mathrm{d}x\right)^{\frac{1}{p}},\quad 0 < p < +\infty.
\end{align*}
用 $L^p(E)$ 表示使 $\|f\|_p < +\infty$ 的 $f$ 的全体，称其为 $\boldsymbol{L}^{\boldsymbol{p}}$ \textbf{空间}.

(ii) 设 $f(x)$ 是 $E\subset\mathbb{R}^n$ 上的可测函数，$m(E) > 0$。若存在 $M$，使得 $|f(x)| \leqslant M$，a.e., $x\in E$，则称 $f(x)$ 在 $E$ 上\textbf{本性有界}，$M$ 称为 $f(x)$ 的\textbf{本性上界}。再对一切本性上界取下确界，记为 $\|f\|_\infty$，称它为 $f(x)$ 在 $E$ 上的本性上确界。此时用 $L^\infty(E)$ 表示在 $E$ 上本性有界的函数之全体。

当 $0 < m(E) < +\infty$ 时，可以证明 $\lim\limits_{p\to\infty} \|f\|_p = \|f\|_\infty$。事实上，不妨记 $M = \|f\|_\infty$，则对任一满足 $M' < M$ 的 $M'$，点集
\begin{align*}
A = \{x\in E: |f(x)| > M'\}
\end{align*}
有正测度。由不等式
\begin{align*}
\|f\|_p = \left(\int_E |f(x)|^p \mathrm{d}x\right)^{\frac{1}{p}} \geqslant M'(m(A))^{\frac{1}{p}}
\end{align*}
可知 $\varliminf\limits_{p\to\infty} \|f\|_p \geqslant M'$。令 $M'\to M$，得 $\varliminf\limits_{p\to\infty} \|f\|_p \geqslant M$。

另一方面，我们总有
\begin{align*}
\|f\|_p \leqslant \left(\int_E M^p \mathrm{d}x\right)^{\frac{1}{p}} = M(m(E))^{\frac{1}{p}},
\end{align*}
从而又得
\begin{align*}
\varlimsup\limits_{p\to\infty} \|f\|_p \leqslant M.
\end{align*}
\end{definition}
\begin{remark}
注意，对 $L^p(E)$，我们主要的兴趣在 $p\geqslant 1$ 的情形，下文中若未指明 $p>0$，则一律认为是 $p\geqslant 1$.
\end{remark}

\begin{theorem}
若 $f,g\in L^p(E)$, $0 < p \leqslant +\infty$, $\alpha,\beta$ 是实数，则
\begin{align*}
\alpha f + \beta g \in L^p(E).
\end{align*}
\end{theorem}

\begin{proof}
(i) 当 $0 < p < \infty$ 时，我们有
\begin{align*}
|\alpha f(x) + \beta g(x)|^p \leqslant 2^p(|\alpha|^p \cdot |f(x)|^p + |\beta|^p \cdot |g(x)|^p).
\end{align*}

(ii) 当 $p = +\infty$ 时，我们有
\begin{align*}
|\alpha f(x) + \beta g(x)| \leqslant |\alpha| \cdot \|f\|_\infty + |\beta| \cdot \|g\|_\infty,\quad \text{a.e. } x\in E,
\end{align*}
从而可知 $\|\alpha f + \beta g\|_\infty \leqslant |\alpha| \cdot \|f\|_\infty + |\beta| \cdot \|g\|_\infty$.

\end{proof}

\begin{example}
设 $E=(0,1)$，则
\begin{gather*}
\ln\frac{1}{x} \in L^p(E),\quad \ln\frac{1}{x} \notin L^\infty(E); \\
x^{-\frac{1}{p}} \in L^{p-\alpha}(E)\quad (0 < \alpha < p),\quad x^{-\frac{1}{p}} \notin L^p(E); \\
x^{-\frac{1}{p}}\left(\ln\frac{1}{x}\right)^{-2/p} \notin L^{p+\alpha}(E)\quad (\alpha > 0), \\
x^{-\frac{1}{p}}\left(\ln\frac{1}{x}\right)^{-2/p} \in L^p(E).
\end{gather*}
\end{example}








\end{document}